%\documentclass{book}\usepackage{sjtfonts}\usepackage{includegraphics}\usepackage{alltt}\usepackage{latexsym}
%\usepackage{array}
%\begin{document}\input{defs0}%		miradefs.tex 		    June 1994

% opening and closing a quoted Miranda section

\newcommand{\mira}{\noindent\hrulefill\nopagebreak\so}
\newcommand{\arim}{\nopagebreak\st\hrulefill\\}

% backslash in the Miranda environment
% Only feasible approach is to use \verb+\+
% in line!

%\newcommand{\bs}{$\mi\backslash\!\!$}
%\newcommand{\lbs}{\mbox{\verb+\+}}

% ignoring part of the input

\newcommand{\ignore}[1]{}

% a blank line in a Miranda section

\newcommand{\bl}{\mbox{\ }}

% Signalling a diffculty.

\definecolor{light-gray}{gray}{0.95}

\newcommand{\beware}[2]
 {\medskip\noindent\colorbox{light-gray}{\parbox{\textwidth}
 {\mbox{\large\textbf{#1}} \medskip\noindent\\ #2}\medskip\noindent}}

% Subscripting in Miranda text
\newcommand{\subscr}[2]{\mbox{${\mbox{\texttt{#1}}}_{\mbox{\texttt{#2}}}$}}

%\newcommand{\subscr}[2]{\mbox{${\mbox{\mi #1}}_{\mbox{\smtt #2}}$}}
\newcommand{\aone}{\subscr{a}{1}}
\newcommand{\atwo}{\subscr{a}{2}}
\newcommand{\ak}{\subscr{a}{k}}

\newcommand{\aoneone}{\subscr{a}{1,1}}

\newcommand{\eone}{\subscr{e}{1}}
\newcommand{\etwo}{\subscr{e}{2}}
\newcommand{\ei}{\subscr{e}{i}}
\newcommand{\er}{\subscr{e}{r}}

\newcommand{\fone}{\subscr{f}{1}}
\newcommand{\fl}{\subscr{f}{l}}

\newcommand{\gone}{\subscr{g}{1}}
\newcommand{\gtwo}{\subscr{g}{2}}
\newcommand{\gi}{\subscr{g}{i}}
\newcommand{\gr}{\subscr{g}{r}}

\newcommand{\hone}{\subscr{h}{1}}
\newcommand{\hl}{\subscr{h}{l}}

\newcommand{\pone}{\subscr{p}{1}}
\newcommand{\ptwo}{\subscr{p}{2}}
\newcommand{\pk}{\subscr{p}{k}}

\newcommand{\qone}{\subscr{q}{1}}
\newcommand{\qtwo}{\subscr{q}{2}}
\newcommand{\qk}{\subscr{q}{k}}

\newcommand{\rone}{\subscr{r}{1}}
\newcommand{\rtwo}{\subscr{r}{2}}

\newcommand{\tone}{\subscr{t}{1}}
\newcommand{\ttwo}{\subscr{t}{2}}
\newcommand{\tn}{\subscr{t}{n}}

\newcommand{\vone}{\subscr{v}{1}}
\newcommand{\vtwo}{\subscr{v}{2}}
\newcommand{\vn}{\subscr{v}{n}}

% for regular expressions

\newcommand{\stone}{\subscr{st}{1}}
\newcommand{\sttwo}{\subscr{st}{2}}
\newcommand{\stn}{\subscr{st}{n}}
\newcommand{\rthree}{\subscr{r}{3}}

\newcommand{\eps}{$\varepsilon$}
\newcommand{\trans}[1]{\stackrel{\mbox{\tt #1}}{\longrightarrow}}



% superscript in Miranda text

\newcommand{\superscr}[2]{\mbox{${\mbox{\mi #1}}^{\mbox{\smtt #2}}$}}

% Not equals in Miranda text

\newcommand{\noteq}{\makebox[0pt][l]{=}/}
\newcommand{\overslash}{\makebox[0pt][l]{/}}

% Definitions in complexity chapter

\newfont{\oldSym}{cmsy10}
\newcommand{\Th}{\boldmath$\oldSym\Theta$}
\newcommand{\pow}[1]{\^{}#1}
\newcommand{\leqq}{$\leq$}
\newcommand{\geqq}{$\geq$}
\newcommand{\lll}{$\ll$}
\newcommand{\eqq}{$\equiv$}
\newcommand{\ltwo}{\subscr{log}{2}}

% Indexing

\newcommand{\minx}[1]{\index{#1@{\ttfamily{}#1}}}
\setcounter{chapter}{19}\setcounter{page}{415}
\addtocontents{toc}{\contentsline {chapter}{Appendices}{}}

%\minitocoptionfalse

\chapter{Functional, imperative and OO programming}
\label{vs}
\index{imperative programming|(}

In this appendix we compare programming in Haskell to more traditional
notions in imperative languages like Pascal and C and object-oriented
(OO)
languages such as C\#, C++ and Java.

\section*{Values and states}
\index{value}\index{state}

Consider the example of finding the sum of squares
of natural numbers up to a particular number. 
A functional program describes the values that are to be calculated, directly. 
\begin{alltt}
sumSquares :: Int -> Int
sumSquares 0 = 0
sumSquares n = n*n + sumSquares (n-1)
\end{alltt}
These equations state what the sum of squares is for a natural number
argument. In the first case it
is a direct description; in the second it states that the sum to non-zero
\texttt{n} is
got by finding the sum to \texttt{n-1} and adding the square of \texttt{n}.

A typical imperative program might solve the problem thus
\begin{alltt}
s = 0 ;
i = 0 ;
while i<n do 
  begin
    i = i+1 ;
    s = i*i + s ;
  end 
\end{alltt}
The sum is the final value of the variable \texttt{s}, which is changed repeatedly
during program execution, as is the `count' variable, \texttt{i}.
The effect of the program can only be seen by following the  sequence of changes
made to these variables by the commands in the program, while  the
functional program can be read as a series of equations defining the sum of squares. This meaning is \textbf{explicit} in the functional program, whereas the imperative program has an overall effect which is not obvious from the program itself. 


The link between these two approaches is given by a tail-recursive solution to the problem in Haskell:
\begin{alltt}
sumSquares n = ssAcc 0 0 n

ssAcc i s n
  | i<n       = ssAcc (i+1) ((i+1)^2+s) n
  | otherwise = s
\end{alltt}
Here the three argument positions play the role of three variables -- \texttt{i}, \texttt{s} and \texttt{n} -- whose values are changed on each call to the loop.

A more striking algorithm still is one which is completely explicit:
`to find the sum of squares, build the list of numbers \texttt{1} to 
\texttt{n}, square each of them, and sum the result'.
This program, which uses neither complex control flow, as does the imperative
example, nor recursion as seen in the function \texttt{sumSquares}, can be written
in a functional style, thus:
\begin{alltt}
newSumSq :: Int -> Int
newSumSq n = sum (map square [1 ..\ n])
\end{alltt}
where \texttt{square x = x*x}, the operation \texttt{map}
applies its first argument to every member of a list, and \texttt{sum} finds the
sum of a list of numbers. More examples of this sort of \textbf{data-directed}
programming\index{data-directed programming}
can be seen in the body of the text.

\section*{Functions and variables}

An important difference between the two styles
is what is meant by some of the terminology. Both `function' and `variable'
have different interpretations.

\index{function}
As was explained earlier, a function in a functional program is simply
something which returns a value which depends upon some inputs. In 
imperative and object-oriented languages like Pascal, C, C++ and Java a 
function is rather different. It will
return a value depending upon its arguments, but in general it will also 
change the
values of variables. Rather than being a pure function it is really
a procedure which returns a value when it terminates. 

\index{variable}
In a functional program a variable stands for an \textbf{arbitrary} or \textbf{
unknown} value. Every occurrence of a variable in an equation is interpreted
in the same way. They are  just like variables in logical formulas,
or the mathematical variables familiar from
equations like
\begin{alltt}
\superscr{a}{2} - \superscr{b}{2} = (a-b)(a+b)
\end{alltt}
In any particular case, the value of all three occurrences of \texttt{a} will be
the same. In exactly the same way, in
\so
sumSquares n = n*n + sumSquares (n-1)
\st
all occurrences of \texttt{n} will be interpreted by the same value. For example
\so
sumSquares 7 = 7*7 + sumSquares (7-1)
\st
The crucial motto is `variables in functional programs {\em do not vary}'.

On the other hand, the value of a variable in an imperative
program changes throughout its lifetime. In the sum of squares program above,
the variable \texttt{s} will take the values \texttt{0,1,5,\ldots} successively.
Variables in imperative programs {\em do\/} vary over time, on the other hand.

\section*{Program verification}
\index{proof}
\vspace{-6pt}

Probably the most important difference between functional and imperative
programs is logical. 
As well as being a program, a functional definition is a logical equation
describing a \textbf{property} of the function. Functional
programs are \textbf{self-describing}, as it were. Using the definitions, other
properties of the functions can be deduced. 

To take a simple example, for all
\texttt{n>0}, it is the case that 
\begin{alltt}
sumSquares n > 0
\end{alltt}
To start with,
\begin{alltt}
sumSquares 1 
= 1*1 + sumSquares 0
= 1*1 + 0
= 1
\end{alltt}
which is greater than 0. In general, for \texttt{n} greater than zero,
\begin{alltt}
sumSquares n = n*n + sumSquares (n-1)
\end{alltt}
Now, \texttt{n*n} is positive, and if \texttt{sumSquares (n-1)} is positive,
their sum, \texttt{sumSquares n}, must be. This proof can be formalized using
\textbf{mathematical induction}\index{mathematical induction}. The body
of the text contains numerous examples of proofs by induction over the
structure of data structures like lists and trees, as well as over
numbers.

Program verification is possible for imperative programs as well, but
imperative programs are not self-describing in the way functional ones are.
To describe the effect of an imperative program, like the `sum of squares'
program above, we need to add to the program logical formulas or assertions
which describe the state of the program at various points in its execution.
These methods are both more indirect and more difficult, and verification
seems very difficult indeed for `real' languages like Pascal and 
C. Another aspect of program verification is \textbf{program transformation} in
which programs are transformed to other programs which have the same effect but
better performance, for example. Again, this is difficult for
traditional imperative languages.


\section*{Records and tuples}
\index{tuples}\index{records}
\vspace{-6pt}

In Chapter \ref{tupleList} the tuple types of Haskell are introduced. In
particular we saw the definition
\begin{alltt}
type Person = (String,String,Int)
\end{alltt}
This compares with a Pascal declaration of a record
\begin{alltt}
type Person = record 
  name  : String;
  phone : String;
  age   : Integer
end;
\end{alltt} 
which has three fields which have to be named. In Haskell the fields of a
tuple can be accessed by pattern matching, but it is possible to define
functions called \textbf{selectors} which behave in a similar way, if
required:  
\begin{alltt}
name  :: Person -> String
name (n,p,a) = n
\end{alltt}
and so on. If \texttt{per ::\ Person} then \texttt{name per ::\ String}, similarly
to  \texttt{r.name} being a string variable if \texttt{r} is a variable of type 
\texttt{Person} in Pascal.

We could instead use an algebraic type to represent a person:
\begin{alltt}
data Person = Person {name::String, phone::String, age::Integer}
              deriving (Show,Eq)
\end{alltt}
In an object-oriented language like Java or C\# this type would be represented by an \emph{object} wit three attributes, one for each of the name, phone and age.  The methods modifying these values would also be part of the definition of the object.

To implement the analogue of an algebraic data type with more than one constructor in Java it is necessary to work rather harder. The type itself is modelled as a class with abstract methods, implemented in a number of sub-classes, one per constructor. Each sub-class contains the attributes for a particular constructor, together with the implementation of the methods over that sub-class. This approach works well for methods that work over one member of a class, but is more problematic for binary methods, and in particular for equality.




\section*{Lists and pointers}
\index{pointer}
\vspace{-5pt}

Haskell contains the type of lists built in, and other recursive types such as
trees can be defined directly. We can think of the type of linked lists given
by pointers in Pascal as an \textbf{implementation} of lists, since in Haskell
it is not necessary to think of pointer values, or of storage allocation (
\texttt{new} and \texttt{dispose}) as it is in Pascal. Indeed, we can think of Haskell
programs as \textbf{designs} for Pascal list programs.
If we define
\index{lists!as Pascal type}
\begin{alltt}
type list = \^{}node;
type node = record
   head : value;
   tail : list
end;
\end{alltt}
then we have the following correspondence, where the Haskell \texttt{head} and
\texttt{tail} functions give the head and tail of a list.
\begin{alltt}
            []                        nil
            head ys                   ys\^{}.head 
            tail ys                   ys\^{}.tail 
            (x:xs)                    cons(x,xs) 
\end{alltt}
The function \texttt{cons} in Pascal has the definition
\begin{alltt}
function cons(y:value;ys:list):list;
  var xs:list;
  begin
    new(xs);
    xs\^{}.head := y;
    xs\^{}.tail := ys;
    cons := xs
  end;
\end{alltt}
Functions such as 
\begin{alltt}
sumList []     = 0 
sumList (x:xs) = x + sumList xs 
\end{alltt}
can then be transferred to Pascal in a
straightforward way.
\begin{alltt}
function sumList(xs:list):integer;
  begin
    if xs=nil 
      then sumList := 0
      else sumList := xs\^{}.head + sumList(xs\^{}.tail)
  end;
\end{alltt}
A second example is 
\begin{alltt}
doubleAll []     = []
doubleAll (x:xs) = (2*x) : doubleAll xs
\end{alltt}
where
we use \texttt{cons} in the Pascal definition of the function
\begin{alltt}
function doubleAll(xs:list):list;
  begin
    if xs=nil 
      then doubleAll := nil
      else doubleAll := cons( 2*xs\^{}.head , doubleAll(xs\^{}.tail) )
  end;
\end{alltt}
If we define the functions
\begin{alltt}
function head(xs:list):value;     function tail(xs:list):list;
  begin                            begin
    head := xs\^{}.head                 tail := xs\^{}.tail
  end;                             end;
\end{alltt}
then the correspondence is even clearer:
\begin{alltt}
function doubleAll(xs:list):list;
  begin
    if xs=nil 
      then doubleAll := nil
      else doubleAll := cons( 2*head(xs) , doubleAll( tail(xs) ) )
  end;
\end{alltt}
This is strong evidence that a functional approach can be useful even if  we
are writing in an imperative language: the functional language can be the
high-level {\em design\/}\index{design}
language for the imperative implementation. Making
this separation can give us substantial help in finding imperative programs
-- we can think about the design and the lower level implementation {\em
separately}, which makes each problem smaller, simpler and therefore easier
to solve.
\vspace{-5pt}

\section*{Higher-order functions}
\index{higher-order function}
\vspace{-5pt}

Traditional imperative languages give little scope for higher-order
programming; Pascal, Java and C allow functions as arguments, so long as those
functions are not themselves higher-order, but has no facility for returning
functions as results. In C++ it is possible to return objects which
represent functions by overloading the function application operator!
This underlies the genericity hailed in the C++ Standard Template Library,
which requires advanced features of the language to implement functions
like \texttt{map} and \texttt{filter}.

Control structures like \texttt{if-then-else} bear some resemblance to
higher-order functions, as they take commands, \texttt{\subscr{c}{1}}, 
\texttt{\subscr{c}{2}} etc.\ into other commands, 
\begin{alltt}
if b then \subscr{c}{1} else \subscr{c}{2}\ \ \ \ \ \ while b do \subscr{c}{1}
\end{alltt}
just as \texttt{map} takes one function to another. Turning the analogy around,
we can think of higher-order functions in Haskell
as \textbf{control structures} which we
can define ourselves. This perhaps explains why we form libraries of
polymorphic functions: they are the control structures we use in programming
particular  sorts of system. Examples in the text include libraries for
building parsers (Section \ref{parsing}) and interactive I/O programs
(Chapter \ref{io}), as well as the built-in list-processing
functions.
\vspace{-2pt}

\section*{Polymorphism}
\index{polymorphism}
\vspace{-2pt}

Again, this aspect is poorly represented in many imperative languages; the
best we can do in \texttt{Pascal}, say, is to use a text editor to copy and
modify the list processing code from one type of lists for use with another.
Of course, we then run the risk that the different versions of the programs
are not modified in step, unless we are very careful to keep track of
modifications, and so on.


Polymorphism in Haskell is what is commonly known as
\textbf{generic} polymorphism: the same `generic' code works over a
whole collection of types. A simple example is the function which 
reverses the elements in a list.

Haskell classes support what is known as `ad
hoc' polymorphism, or in object-oriented terminology simply `polymorphism', 
in which different programs implement the same
operation over different types. An example of this is the \texttt{Eq}
class of types carrying an equality operation: the way in which equality
is checked is completely different at different types. Another way of
viewing classes is as \textbf{interfaces}\index{interface} which different types
can implement in different ways; in this way they resemble the
interfaces of object-oriented languages like Java.

As is argued in the text, polymorphism is one of the mechanisms which helps to
make programs {\em reusable\/} in Haskell; it remains to be seen whether this
will also be true of advanced imperative languages.
\vspace{-2pt}

\section*{Defining types and classes}
\vspace{-2pt}

The algebraic type mechanism of Haskell, explained in Chapter \ref{algTypes},
\index{algebraic type}
subsumes various traditional type definitions. Enumerated types are given by
algebraic types all of whose constructors are 0-ary (take no arguments);
variant records can be implemented as algebraic types with more then one
constructor, and \textbf{recursive} types usually implemented by means of
pointers become recursive algebraic types.

Just as we explained for lists, Haskell programs over trees and so on can be
seen as {\em designs\/} for programs in imperative languages manipulating the
pointer implementations of the types.

The abstract data types\index{abstract data type},
introduced in Chapter \ref{adt}, are very like the abstract
data types of \texttt{Modula-2} and so on; the design  methods we suggest for
use of abstract data types mirror
aspects of the \textbf{object-based} approach advocated for
modern imperative languages such as \texttt{Ada}. 

\index{classes}
The Haskell class system also has object-oriented aspects, as we saw in
Section \ref{algClasses}. It is important to note that Haskell classes are in
some ways quite different from the classes of, for instance, \texttt{C++}. In
Haskell classes are made up of types, which themselves have members; in 
\texttt{C++} a class is like a type, in that it contains objects. Because of this many
of the aspects of object-oriented design in \texttt{C++} are seen as issues of
type design in Haskell.
\vspace{-3pt}

\section*{List comprehensions}
\index{list comprehensions}
\vspace{-3pt}

List comprehensions provide a convenient notation for \textbf{iteration\/} along
lists: the analogue of a \texttt{for} loop, which can be used to run through the
indices of an array. For instance, to sum all pairs of elements of \texttt{xs} and
\texttt{ys}, we write
\begin{alltt}
[ a+b | a <- xs , b <- ys ]
\end{alltt}
The order of the iteration is for a value \texttt{a} from the list \texttt{xs} to be
fixed and then for \texttt{b} to run through the possible values from \texttt{ys};
this is then repeated with the next value from \texttt{xs}, until the list is
exhausted. Just the same happens for a \textbf{nested\/} \texttt{for} loop
\begin{alltt}
for i:=0 to xLen-1 do
  for j:=0 to yLen-1 do\hfill (twoFor)
    write( x[i]+y[j] ) 
\end{alltt}
where we fix a value for \texttt{i} while running through all values for 
\texttt{j}.

In the \texttt{for} loop, we have to run through the indices; a list generator
runs through the values directly. The indices of the list
\texttt{xs} are given by 
\begin{alltt}
[0 ..\ length xs - 1]
\end{alltt}
and so a Haskell analogue of 
\texttt{(twoFor)} can be written thus:
\begin{alltt}
[ xs!!i + ys!!j | i <- [0 ..\ length xs - 1] , 
                  j <- [0 ..\ length ys - 1] ]
\end{alltt}
if we so wish. 
\vspace{-3pt}

\section*{Lazy evaluation}
\index{lazy evaluation}
\vspace{-3pt}

Lazy evaluation and imperative languages do not mix well. In \texttt{Pascal},
for instance, we can write the function definition
\begin{alltt}
function succ(x : integer):integer;
begin
  y    := y+1;
  succ := x+1
end;
\end{alltt}
This function adds one to its argument, but also has the \textbf{side-effect\/}
of increasing \texttt{y} by one. If we evaluate \texttt{f(y,succ(z))} we cannot
predict the effect it will have.
\begin{titemize}
\item If \texttt{f} evaluates its second argument first, \texttt{y} will be
increased before being passed to \texttt{f}; 
\item on the other hand, if \texttt{f} needs its first argument first (and perhaps its second
argument not at all), the value passed to \texttt{f} will not be increased, even
if it is increased before the function call terminates.
\end{titemize}
In general, it will not be possible to predict the behaviour of even the
simplest programs. Since evaluating an expression can cause a change of the
state, the order of expression evaluation determines the overall effect of a
program, and so a lazy implementation can behave differently (in unforeseen
ways) from the norm.

\section*{State, infinite lists and monads}
\index{lists!infinite}\index{state}\index{monad}

Section \ref{infLists} introduced infinite lists, and one of the first
examples given there was an infinite list of random numbers. This list could
be supplied to a function requiring a supply of random numbers; because of lazy
evaluation, these numbers will only be generated on demand. 

If we were to implement this imperatively, we would probably keep in a
variable the last random number generated, and at each request for a
number we would update this store. We can see the infinite list as supplying
{\em all the values that the variable will take\/} as a single structure; we
therefore do not need to keep the state, and hence have an \textbf{abstraction\/} from the imperative view.

We have seen in Section \ref{monadFP} that there has been recent important
work on integrating side-effecting programs into a functional system by a
monadic approach.

\section*{Conclusion}

Clearly there are parallels between the functional and the imperative, as
well as clear differences. The functional view of a system is often
higher-level, and so even if we ultimately aim for an imperative solution, a
functional design  or \textbf{prototype\/} can be most useful.

We have seen that monads can be used to give an interface to 
imperative features within a functional framework. Many of the Haskell
implementations offer these facilities, and so give a method of uniting the
best features of two important programming paradigms without compromising the
purity of the language. Other languages,
including Standard ML \cite{defsml} and F\# \cite{Fsharp}, combine the functional and the imperative, but
these systems tend to lose their pure functional properties in the process.

It is interesting to see the influence of ideas from modern functional
programming languages in the design of Java extensions. One of the main
drawbacks of Java for a long time was that it lacked generic polymorphism. The current 
mechanism for generics in the Java standard owes its inspiration and much of its detail to Haskell polymorphism.


\index{imperative programming|)}

%\end{document}
